\documentclass[10pt,a4paper]{amsart}
\usepackage[margin=19mm]{geometry}
\usepackage{amsmath,amssymb,mathtools}
\usepackage[scr=rsfs,cal=euler]{mathalfa}
\usepackage{xfrac}
\usepackage[unicode,hidelinks,bookmarks=false]{hyperref}
\newcommand{\ie}{i.e.\ }
\newcommand{\cf}{cf.\ }
\newcommand{\resp}{respectively\ }
\newcommand{\Subsection}{\subsection}
\providecommand{\nobreakdash}{\nobreak\hbox{-}}
\setlength{\emergencystretch}{2em}
\multlinegap0pt
\usepackage{bm}
\usepackage{bbm}
\usepackage{dsfont}
\usepackage{xspace}
\usepackage{cancel}
\usepackage{mathtools}
\usepackage{amsthm}
\makeatletter
\@ifundefined{theorem}{\newtheorem{theorem}{Theorem}}{}
\@ifundefined{definition}{\newtheorem{definition}{Definition}}{}
\@ifundefined{lemma}{\newtheorem{lemma}[theorem]{Lemma}}{}
\makeatother
\title[Competing Bandits in Matching Markets via Super Stability]{Competing Bandits in Matching Markets via Super Stability}
\author{Soumya Basu}
\date{Source: arXiv:2506.15926v1 (CC BY 4.0)}
\begin{document}
\maketitle

\noindent\emph{Source and licence.} Competing Bandits in Matching Markets via Super Stability. Version arXiv:2506.15926v1 (CC BY 4.0), distributed under the Creative Commons Attribution 4.0 International licence, as recorded in the arXiv licence metadata for this exact version and shown on the abstract page https://arxiv.org/abs/2506.15926v1. Full rights notice: licenses/arxiv-2506.15926-CC-BY-4.0.txt.\par\smallskip

\noindent\textit{Editorial scope.} Counted unit: the Lemma whose source label is \texttt{lemm:opt\_lambda\_c} in the article (source0.tex lines 1247--1249, source label \texttt{lemm:opt\_lambda\_c}), delivered together with the article's own complete original proof of it (source0.tex lines 1300--1337, one \texttt{proof} environment of 38 lines). No mathematics is written, repaired or re-derived: statement and proof are the article's own text line for line. The proof is detached from the statement in the source: the article states the result at lines 1247--1249 and prints its proof at lines 1300--1337, and the proof environment's own header names the statement, so the association is the article's own. Required context retained uncounted: def:cover (lines 1071--1085); eq:lower\_4 (lines 1076--1084); eq:lower\_5 (lines 1076--1084); eq:lower\_6 (lines 1076--1084); eq:kl\_theta\_c (lines 1095--1097). Every definition, display and label of those lines is retained unchanged. Omitted from this package and not redistributed: 1--1070, 1085--1094, 1098--1246, 1250--1299, 1338--1340. Nothing inside a retained excerpt is omitted or altered: every retained body is delivered exactly as the article's lines are written, so no omission receipt of any kind accompanies this package. Citations: the retained excerpts carry no citation command at all, so no bibliography entry of the article is transcribed and no reference is redistributed. Reference adaptation: none was needed and none was made; every reference inside the delivered unit resolves inside it, so no unresolved reference or citation is produced. Printed slips kept literal: no printed word, symbol, hypothesis or number of the counted statement or of its proof was altered. Layout and class: the article's own class is \texttt{article} with packages including PRIMEarxiv, inputenc, fontenc, hyperref, url, booktabs, amsfonts, nicefrac, microtype, lipsum, fancyhdr, graphicx, bbm, enumitem; the wrapper is the lane's pinned \texttt{amsart} editorial wrapper at 10pt on A4, which restates the theorem-like environments the retained text needs and adds the article's own macro definitions that the retained text needs (none), with extra wrapper packages (bm, bbm, dsfont, xspace, cancel, mathtools). The delivered numbering is the unit's own. Assets: no retained span contains a figure environment, a graphics-inclusion command, a PDF-page inclusion, a table or a TikZ picture; no image file is copied into this package, the package declares no asset dependency and no image file is redistributed with it. Licence: version arXiv:2506.15926v1 of this article is distributed under the Creative Commons Attribution 4.0 International licence, as recorded in the arXiv licence metadata for this exact version and shown on the abstract page https://arxiv.org/abs/2506.15926v1. A full rights notice is delivered at licenses/arxiv-2506.15926-CC-BY-4.0.txt. Characters: every retained excerpt is pure ASCII, so the delivered engine, XeLaTeX with its default Latin Modern text font, reports no missing character.
\begin{definition}[Cover of $\mathcal{B}(\theta)$] \label{def:cover}
A set of triplets $(i,j,j')$ or $(j,i,i')$ for $i,i'\in [N]$ and $j,j'\in [K]$, $C$ is a minimal cover of $\mathcal{B}(\theta)$, i.e. $C \in Cover(\theta)$, \emph{iff} all the following statements hold\\
(i) each $(P_u, P_a) \in \mathcal{B}(\theta)$ there exists a triplet $o \in C$ such that $o$ is  $\theta$-open for $(P_u, P_a)$, \\ 
(ii) any proper subset $C' \subset C$, $C'$ is not a cover of $\mathcal{B}(\theta)$,\\
(iii) the set $\Lambda_C$ defined below is non empty
\begin{align}
(\boldsymbol{\tilde{\mu}}^{C}, \boldsymbol{\tilde{\gamma}}^{C}) \in \Lambda_C   \iff &(\boldsymbol{\tilde{\mu}}^{C}, \boldsymbol{\tilde{\gamma}}^{C}) \text{ satisfies} \notag\\
%%%
&\tilde{\mu}_{i,j}^{C} =  \mu_{i,j}(\theta), \tilde{\gamma}_{j,i}^{C} =  \gamma_{j,i}(\theta), \quad\forall (i,j) \in Locked(\theta),  \label{eq:lower_4}\\
%%%
&\tilde{\mu}_{i,j'}^{C} \geq \tilde{\mu}_{i,j}^{C} + \varepsilon,  \quad\forall (i, j, j') \in C,  \label{eq:lower_5}\\
%%%
&\tilde{\gamma}_{j,i'}^{C} \geq \tilde{\gamma}_{j,i}^{C} + \varepsilon, \quad\forall (j, i, i') \in C, \label{eq:lower_6}
\end{align}
\end{definition} 

\begin{equation}\label{eq:kl_theta_c}
 kl(\theta, C) \triangleq\sum_{i}\max_{\substack{(i,j) \notin Locked(\theta),\\ (i,j')\in Locked(\theta)\\ (i,j,j'), (i,j',j) \in C}} kl\big( \mu_{ij}(\theta), \mu_{ij'}(\theta) \big) + \sum_{j}\max_{\substack{(i,j) \notin Locked(\theta),\\ (i',j)\in Locked(\theta)\\ (j,i,i'), (j,i',i) \in C}} kl\big(\gamma_{ji}(\theta), \gamma_{ji'}(\theta)\big),
\end{equation}

\begin{lemma}\label{lemm:opt_lambda_c}
For a minimal cover $C \in Cover(\theta)$ there exists a $\lambda_C$ that for all $M \notin Stable(\theta)$  satisfies $kl(\theta, \lambda_C; M) \leq kl(\theta, C)$ where $kl(\theta, C)$ is as defined in Equation~\eqref{eq:kl_theta_c}.
\end{lemma}

\begin{proof}[Proof of Lemma~\ref{lemm:opt_lambda_c}]
We fix an arbitrary cover $C \in Cover(\theta)$. First, observe that for any $C \in Cover(\theta)$ we first note that keeping  $\tilde{\mu}^C_{i,j} = \mu_{i,j}(\theta)$ for all $(i,j)$ such that $(i, \cdot, j) \notin C$  and $(i, j, \cdot) \notin C$ satisfies the inequalities \eqref{eq:lower_5}. Similarly,  keeping $\tilde{\gamma}^C_{j,i} = \gamma_{j,i}(\theta)$ for all $(j,i)$ such that $(j, \cdot, i) \notin C$  and $(j, i, \cdot) \notin C$ satisfies the inequalities \eqref{eq:lower_6}. So the changes are isolated in the triplets under $C$. Therefore, it follows that 
\begin{align*}
&kl(\theta, \lambda_C; M) \\
&= \sum_{(i,j)\in M}kl(\mu_{i,j}(\theta), \tilde{\mu}^C_{i,j}) + kl(\gamma_{j,i}(\theta), \tilde{\gamma}^C_{j,i})\\
%
&\leq \sum_{i}\max_{j:(i,j)\notin Locked(\theta)} kl(\mu_{i,j}(\theta), \tilde{\mu}^C_{i,j}) + \sum_{j}\max_{i:(i,j)\notin Locked(\theta)} kl(\gamma_{j,i}(\theta), \tilde{\gamma}^C_{j,i})\\
%
&\leq \sum_{i}\max_{j:(i,j)\notin Locked(\theta)}\left(\mathbbm{1}((i,\cdot,j)\in C\, \vee (i,j, \cdot) \in C) kl(\mu_{i,j}(\theta), \tilde{\mu}^C_{i,j}) \right) \\
&+  \sum_{j}\max_{i:(i,j)\notin Locked(\theta)} \left(\mathbbm{1}((j,\cdot,i)\in C\, \vee (j, i, \cdot) \in C) kl(\gamma_{j,i}(\theta), \tilde{\gamma}^C_{j,i})\right)\\
%
&\leq \sum_{i}\max_{j:(i,j)\notin Locked(\theta)}\max_{\delta = \pm\varepsilon}\max_{\substack{(i,j')\in Locked(\theta)\\ (i,j,j'), (i, j', j)\in C}} kl(\mu_{i,j}(\theta), \mu_{i,j'}(\theta) + \delta) \\
&+ \sum_{j}\max_{i:(i,j)\notin Locked(\theta)} 
\max_{\delta = \pm\varepsilon}\max_{\substack{(i',j)\in Locked(\theta)\\ (j,i,i'), (j, i', i)\in C}} kl(\gamma_{i,j}(\theta), \gamma_{i',j}(\theta) + \delta)
\end{align*}
The first inequality is true as all non-zero $kl$ distance values are isolated to $(i,j) \notin Locked(\theta)$ as per \eqref{eq:lower_4}. The second inequality holds as the changes are further isolated to triplets inside the cover $C$. We now explain the third inequality. 
For a given cover $C$, by Definition~\ref{def:cover}  there is a valid $\lambda_C$ that satisfies equations \eqref{eq:lower_4}, \eqref{eq:lower_5}, and \eqref{eq:lower_6}. For such a $\lambda_C$, we  claim to have for any $(i,j) \notin Locked(\theta)$
\begin{align}
    & \min_{\substack{(i,j')\in Locked(\theta)\\(i, j, j') \in C}} \mu_{i,j'}(\theta) -\varepsilon \leq \tilde{\mu}^C_{i,j}  \leq   \max_{\substack{(i,j')\in Locked(\theta)\\(i, j', j) \in C}} \mu_{i,j'}(\theta) + \varepsilon,\label{eq:mu_ineq}\\
    &\min_{\substack{(i,j')\in Locked(\theta)\\(j, i, i') \in C}} \gamma_{j,i'}(\theta) -\varepsilon \leq \tilde{\gamma}^C_{j,i} \leq \max_{\substack{(i,j')\in Locked(\theta)\\(j, i', i) \in C}} \gamma_{j,i'}(\theta) + \varepsilon,\label{eq:gamma_ineq}
\end{align}
Due to the inequalities \eqref{eq:lower_4} and \eqref{eq:lower_5} we have 
\begin{align*}
&\tilde{\mu}^C_{i,j} \leq \mu_{i,j'}(\theta) -\varepsilon, \forall j': (i, j, j') \in C, (i,j') \in Locked(\theta), \\
&\tilde{\mu}^C_{i,j} \geq \mu_{i,j'}(\theta) + \varepsilon, \forall j': (i, j', j) \in C,  (i,j') \in Locked(\theta).
\end{align*}
Therefore Equation~\eqref{eq:mu_ineq} holds. The inequality \eqref{eq:gamma_ineq} follows in a similar manner.

It is easy to see  the inequality ~\eqref{eq:mu_ineq} implies 
$$
kl(\mu_{i,j}(\theta), \tilde{\mu}^C_{i,j}) \leq \max_{\delta = \pm\varepsilon}\max_{\substack{(i,j')\in Locked(\theta)\\ (i,j,j'), (i, j', j)\in C}} kl(\mu_{i,j}(\theta), \mu_{i,j'}(\theta) + \delta).
$$
Similarly, the inequality ~\eqref{eq:gamma_ineq} implies
$$
kl(\gamma_{j,i}(\theta), \tilde{\gamma}^C_{i,j}) \leq \max_{\delta = \pm\varepsilon}\max_{\substack{(i',j)\in Locked(\theta)\\(j,i,i'), (j, i', i)\in C}} kl(\gamma_{i,j}(\theta), \gamma_{i',j}(\theta) + \delta).
$$
Finally, taking infimum over $\varepsilon>0$ we obtain the final value of $kl(\theta, C)$ in Equation~\ref{eq:kl_theta_c}.
\end{proof}

\end{document}
